% ============================================================
%  MTH636 – Game Theory | Lecture 16
%  Extensive Form Games & Perfect Information
% ============================================================
\documentclass[12pt, a4paper]{article}

% ============================================================
%  MTH636 -- Game Theory
%  Notes Preamble — Color-Coded, TikZ-Enhanced
% ============================================================

% --- Core packages ---
\usepackage[T1]{fontenc}
\usepackage[utf8]{inputenc}
\usepackage{lmodern}
\usepackage[margin=2.5cm, top=2.8cm, bottom=2.8cm]{geometry}
\usepackage{microtype}

% --- Math ---
\usepackage{amsmath, amssymb, amsthm, mathtools}
\usepackage{bm}        % bold math
\usepackage{cancel}    % strike-through in math

% --- Color & Boxes ---
\usepackage[dvipsnames, svgnames, x11names, table]{xcolor}
\usepackage[most]{tcolorbox}
\tcbuselibrary{breakable, skins, theorems}

% --- TikZ & Diagrams ---
\usepackage{tikz}
\usepackage{tikz-qtree}
\usetikzlibrary{arrows.meta, positioning, shapes.geometric, fit, backgrounds,
                decorations.pathreplacing, calc, matrix, trees}
\usepackage{pgfplots}
\pgfplotsset{compat=1.18}

% --- Tables ---
\usepackage{booktabs}
\usepackage{multirow}
\usepackage{array}
\usepackage{tabularx}
\usepackage{makecell}

% --- Typography & Layout ---
\usepackage{enumitem}
\usepackage{parskip}
\usepackage{fancyhdr}
\usepackage{titlesec}
\usepackage{titletoc}
\usepackage{hyperref}
\usepackage{graphicx}
\usepackage{float}
\usepackage{caption}

% --- Color Palette ---
\definecolor{defblue}{HTML}{1A5276}      % Definition header
\definecolor{defbluebg}{HTML}{D6EAF8}    % Definition background
\definecolor{thmgreen}{HTML}{1E8449}     % Theorem header
\definecolor{thmbg}{HTML}{D5F5E3}        % Theorem background
\definecolor{exorange}{HTML}{B7770D}     % Example header
\definecolor{exbg}{HTML}{FEF9E7}         % Example background
\definecolor{warnred}{HTML}{922B21}      % Warning/Key Insight header
\definecolor{warnbg}{HTML}{FADBD8}       % Warning background
\definecolor{noticegray}{HTML}{2C3E50}   % Note header
\definecolor{noticebg}{HTML}{EBF5FB}     % Note background
\definecolor{taketawayyellow}{HTML}{7D6608} % Key Takeaway header
\definecolor{takeawaybg}{HTML}{FDFEFE}   % Takeaway background
\definecolor{accentpurple}{HTML}{6C3483} % Accent / proof
\definecolor{lightgray}{HTML}{F2F3F4}    % shading

% ============================================================
%  TCOLORBOX ENVIRONMENTS
% ============================================================

% -- Definition (Blue) --
\newtcolorbox{defbox}[2][]{
  enhanced, breakable,
  colback=defbluebg, colframe=defblue, coltitle=white,
  fonttitle=\bfseries\sffamily,
  title={Definition: #2},
  attach boxed title to top left={yshift=-2mm, xshift=4mm},
  boxed title style={colback=defblue, rounded corners},
  arc=4pt, boxrule=1.2pt,
  left=6pt, right=6pt, top=8pt, bottom=6pt,
  #1
}

% -- Theorem / Proposition (Green) --
\newtcolorbox{thmbox}[2][]{
  enhanced, breakable,
  colback=thmbg, colframe=thmgreen, coltitle=white,
  fonttitle=\bfseries\sffamily,
  title={Theorem: #2},
  attach boxed title to top left={yshift=-2mm, xshift=4mm},
  boxed title style={colback=thmgreen, rounded corners},
  arc=4pt, boxrule=1.2pt,
  left=6pt, right=6pt, top=8pt, bottom=6pt,
  #1
}

% -- Proposition (Green variant) --
\newtcolorbox{propbox}[2][]{
  enhanced, breakable,
  colback=thmbg, colframe=thmgreen!70!black, coltitle=white,
  fonttitle=\bfseries\sffamily,
  title={Proposition: #2},
  attach boxed title to top left={yshift=-2mm, xshift=4mm},
  boxed title style={colback=thmgreen!70!black, rounded corners},
  arc=4pt, boxrule=1.2pt,
  left=6pt, right=6pt, top=8pt, bottom=6pt,
  #1
}

% -- Example (Orange/Yellow) --
\newtcolorbox{exbox}[2][]{
  enhanced, breakable,
  colback=exbg, colframe=exorange, coltitle=white,
  fonttitle=\bfseries\sffamily,
  title={Example: #2},
  attach boxed title to top left={yshift=-2mm, xshift=4mm},
  boxed title style={colback=exorange, rounded corners},
  arc=4pt, boxrule=1.2pt,
  left=6pt, right=6pt, top=8pt, bottom=6pt,
  #1
}

% -- Key Insight / Warning (Red) --
\newtcolorbox{warnbox}[2][]{
  enhanced, breakable,
  colback=warnbg, colframe=warnred, coltitle=white,
  fonttitle=\bfseries\sffamily,
  title={Key Insight: #2},
  attach boxed title to top left={yshift=-2mm, xshift=4mm},
  boxed title style={colback=warnred, rounded corners},
  arc=4pt, boxrule=1.2pt,
  left=6pt, right=6pt, top=8pt, bottom=6pt,
  #1
}

% -- Note (Gray-Blue) --
\newtcolorbox{notebox}[1][]{
  enhanced, breakable,
  colback=noticebg, colframe=noticegray!60, coltitle=white,
  fonttitle=\bfseries\sffamily,
  title={Note},
  attach boxed title to top left={yshift=-2mm, xshift=4mm},
  boxed title style={colback=noticegray!60, rounded corners},
  arc=4pt, boxrule=0.8pt,
  left=6pt, right=6pt, top=8pt, bottom=6pt,
  #1
}

% -- Key Takeaway --
\newtcolorbox{takebox}[1][]{
  enhanced, breakable,
  colback=takeawaybg, colframe=taketawayyellow, coltitle=white,
  fonttitle=\bfseries\sffamily,
  title={Key Takeaway},
  attach boxed title to top left={yshift=-2mm, xshift=4mm},
  boxed title style={colback=taketawayyellow, rounded corners},
  arc=4pt, boxrule=1.2pt,
  left=6pt, right=6pt, top=8pt, bottom=6pt,
  #1
}

% -- Intuition --
\newtcolorbox{intubox}[1][]{
  enhanced, breakable,
  colback=accentpurple!8, colframe=accentpurple!60, coltitle=white,
  fonttitle=\bfseries\sffamily,
  title={Intuition},
  attach boxed title to top left={yshift=-2mm, xshift=4mm},
  boxed title style={colback=accentpurple!80, rounded corners},
  arc=4pt, boxrule=0.8pt,
  left=6pt, right=6pt, top=8pt, bottom=6pt,
  #1
}

% ============================================================
%  PAYOFF MATRIX HELPER
%  Usage: inside tabular, use \payoff{a}{b} for (a,b)
% ============================================================
\newcommand{\payoff}[2]{$(#1,\;#2)$}
\newcommand{\nash}[2]{\underline{$#1$},\underline{$#2$}}  % underline NE payoffs
\newcommand{\Nset}{\mathbb{N}}
\newcommand{\Rset}{\mathbb{R}}
\newcommand{\Sset}{\mathcal{S}}
\newcommand{\Ui}{u_i}

% ============================================================
%  SECTION FORMATTING
% ============================================================
\titleformat{\section}
  {\large\bfseries\sffamily\color{defblue}}
  {\thesection}{1em}{}[\titlerule]

\titleformat{\subsection}
  {\normalsize\bfseries\sffamily\color{thmgreen!80!black}}
  {\thesubsection}{1em}{}

\titleformat{\subsubsection}
  {\normalsize\itshape\sffamily\color{exorange!80!black}}
  {\thesubsubsection}{1em}{}

% ============================================================
%  HEADER / FOOTER
% ============================================================
\pagestyle{fancy}
\fancyhf{}
\fancyhead[L]{\sffamily\small\color{defblue} MTH636 -- Game Theory}
\fancyhead[R]{\sffamily\small\color{defblue} \leftmark}
\fancyfoot[C]{\sffamily\small\thepage}
\renewcommand{\headrulewidth}{0.4pt}

% ============================================================
%  HYPERREF SETUP
% ============================================================
\hypersetup{
  colorlinks=true,
  linkcolor=defblue,
  urlcolor=thmgreen,
  citecolor=accentpurple,
  bookmarksnumbered=true,
  pdftitle={MTH636 Game Theory Notes},
  pdfauthor={MTH636 Student Notes}
}

% ============================================================
%  MISC HELPERS
% ============================================================
\newcommand{\BR}{\text{BR}}        % Best Response
\newcommand{\NE}{\text{NE}}        % Nash Equilibrium
\newcommand{\SPNE}{\text{SPNE}}    % Subgame Perfect NE
\newcommand{\Dom}{\text{dom}}      % Dominance
\newcommand{\mixs}[1]{\sigma_{#1}} % Mixed strategy
\newcommand{\pures}[1]{s_{#1}}     % Pure strategy


% ============================================================
\begin{document}

% Title
\begin{center}
  {\LARGE\bfseries\sffamily\color{defblue} MTH636 -- Game Theory}\\[6pt]
  {\Large\sffamily Lecture 16: Extensive Form Games \& Perfect Information}\\[4pt]
\end{center}
\vspace{0.5cm}
\hrule height 1.5pt
\vspace{0.3cm}

% ============================================================
\section{Directed Rooted Trees}

To formally define an extensive form game, we first need to define the underlying graph structure.

\begin{defbox}{Directed Rooted Tree}
  A \textbf{directed rooted tree} is a triplet $T = (V, E, n_0)$ where:
  \begin{itemize}
      \item $V$ is a finite set of elements called \textbf{nodes} or vertices (with $|V| \geq 2$).
      \item $E \subseteq V \times V$ are the \textbf{edges}. An edge $e = (n,y) \in E$ is an outgoing edge of $n$ and an incoming edge of $y$.
      \item $n_0 \in V$ is called the \textbf{root} of the tree. For every $n \in V \setminus \{n_0\}$, there is a unique path from $n_0$ to $n$. The root $n_0$ has no incoming edges.
  \end{itemize}
\end{defbox}

\begin{itemize}
    \item \textbf{Leaves ($Z$):} Nodes that do not have any outgoing edges. Denoted by $Z \subseteq V$.
    \item \textbf{Children ($C(n)$):} For a node $n \in V$, the children are $C(n) = \{y \in V \mid (n,y) \in E\}$.
    \item \textbf{Actions ($A(n)$):} For a node $n \in V$, the available actions are the outgoing edges $A(n) = \{e \in E \mid e = (n,y) \text{ for some } y \in V\}$.
\end{itemize}

% ============================================================
\section{Formal Definition of Extensive Form Games}

\begin{defbox}{Extensive Form Game (Perfect Information)}
  An extensive form game is a tuple $\Gamma = \langle N, V, E, n_0, (V_i)_{i \in N}, (u_i)_{i \in N} \rangle$ where:
  \begin{itemize}
      \item $N$ is a finite set of players.
      \item $(V, E, n_0)$ is a directed rooted tree.
      \item $(V_i)_{i \in N}$ is a partition of the non-leaf nodes (the nodes where decisions are made), i.e., a partition of the set $V \setminus Z$.
      \item $(u_i)_{i \in N}$ is a collection of utility functions, where $u_i : Z \to \mathbb{R}$ assigns a payoff to every leaf node for player $i$.
  \end{itemize}
\end{defbox}

Let $J: V \setminus Z \to N$ be a function such that $J(n) = i$ if $n \in V_i$. 
The game proceeds in the following way:
\begin{enumerate}
    \item Player $J(n_0)$ initiates the game by choosing a possible action from $A(n_0)$. Equivalently, he chooses a node $n_1 \in C(n_0)$.
    \item If $n_1$ is a leaf, the game ends there. Otherwise, player $J(n_1)$ chooses a node $n_2 \in C(n_1)$.
    \item In this way, the game continues until a leaf node is reached, at which point payoffs are distributed according to $(u_i)$.
\end{enumerate}

\begin{notebox}
\textbf{Perfect Information:}
These games are perfect information games. The players know exactly what has happened before making their moves (they know exactly which node they are at).
\end{notebox}

% ============================================================
\section{Strategies in Extensive Form Games}

\begin{defbox}{Pure Strategy}
  A pure strategy for player $i$, denoted $s_i$, is a mapping which maps each node $n \in V_i$ (the nodes player $i$ controls) to an element of $A(n)$ (an available action at that node). 
  
  Equivalently, it maps $n$ to a child node in $C(n)$.
  \[ s_i : V_i \to \bigcup_{n \in V_i} A(n) \]
  such that $s_i(n) \in A(n)$ for all $n \in V_i$.
\end{defbox}

\textbf{Crucial Note:} A strategy must define an action for \textit{every} node in $V_i$, even for nodes that might never be reached if the player follows that very strategy earlier in the tree. This is necessary to evaluate deviations and subgames.

\begin{exbox}{Strategy Example}
  If Player 1 controls nodes $n_0$ and $n_4$, a valid strategy $s_1$ must specify a choice at both nodes. For example:
  \[ s_1(n_0) = R \quad \text{and} \quad s_1(n_4) = A \]
  This strategy is denoted as $(R, A)$. Even if playing $R$ at $n_0$ means node $n_4$ is never reached, the strategy still requires specifying action $A$ at node $n_4$.
\end{exbox}

\end{document}
